\section{Background}

In this Section, we provide background on existing scenario datasets and the ScTrans dataset that we use as the baseline in our adversarial generation process. We also briefly discuss the CARLA simulation platform and the two ADS used in our approach and evaluation. Then, we introduce the NHTSA pre-crash typology that categorises vehicle movements and interactions immediately prior to a crash. Finally, we discuss the thresholds that we use to classify scenarios as safety-critical based on the Euro NCAP Assisted Driving Protocol. 
%We first review existing work in scenario-based ADS testing and adversarial scenario generation, identifying the gap that motivates our approach (Sec.~\ref{sec:scenario_based_testing}--\ref{sec:scenario_generation}). We then introduce the simulation platforms and ADS used in our project (Sec.~\ref{sec:platform_ads}), followed by the two external frameworks --- the NHTSA pre-crash scenario typology and the Euro NCAP assisted driving protocol --- that ground our methodology in real-world crash statistics and regulatory kinematic standards (Sec.~\ref{sec:NHTSA}--\ref{sec:euro_ncap}).
\subsection{Scenario Datasets for testing ADS}\label{sec:scenario_based_testing}

Simulation-based testing has become the dominant paradigm for ADS evaluation, supported by naturalistic driving datasets such as inD~\cite{bock_ind_2020}, highD~\cite{krajewski_highd_2018}, rounD~\cite{krajewski_round_2020}, Waymo Open Dataset~\cite{sun_scalability_2020}, and nuScenes~\cite{caesar_nuscenes_2020}, as well as the CARLA Leaderboard~\cite{noauthor_carla_nodate} benchmark derived from the NHTSA pre-crash typology. We use the SCTrans dataset~\cite{dai_sctrans_2024} as the basis for our seed scenario identification. SCTrans constructs over 1,900 simulation scenarios by transforming naturalistic traffic recordings into OpenSCENARIO files compatible with CARLA and LGSVL, and demonstrates that its scenarios can boost the collision-finding capability of existing ADS fuzzers by approximately seven times. However, a fundamental limitation shared by all these datasets is that they model \textit{well-behaved} traffic participants whose actions conform to expected driving norms, leaving adversarial edge cases underrepresented.

\subsection{Simulation Platform and ADS Under Test}\label{sec:platform_ads}

We use the CARLA simulator~\cite{dosovitskiy_carla_2017}, an open-source driving platform built on Unreal Engine that provides high-fidelity rendering, realistic vehicle physics, and flexible sensor models. It supports the OpenSCENARIO and OpenDRIVE standards for reproducible scenario execution and has been widely adopted in ADS testing research~\cite{kim_drivefuzz_2022, yan_-demand_2025, li_av-fuzzer_2020}. Two ADS are evaluated:

\textbf{Apollo}~\cite{noauthor_apollo_nodate} is an open-source, industrial-grade L4 ADS integrating perception (LiDAR and camera fusion), prediction, planning, and control into a full-stack, machine-learning-based pipeline. It has been extensively used as a test subject in ADS research~\cite{dai_sctrans_2024, kim_drivefuzz_2022, yan_-demand_2025} due to its full-stack architecture and real-world deployment.

\textbf{CARLA Traffic Manager (TM)}~\cite{dosovitskiy_carla_2017} is the built-in rule-based autopilot that controls vehicle behaviour through parameterizable rules governing target speed, following distance, lane-change probability, and traffic light compliance. We include it to compare the robustness of fundamentally different ADS architectures --- one machine-learning-based and one rule-based --- against identical adversarial scenarios.

We use 854 successfully simulated scenarios from the SCTrans dataset~\cite{dai_sctrans_2024} as the starting point for our adversarial generation pipeline. SCTrans formalises the conversion from naturalistic traffic recordings to OpenSCENARIO files as a Model Transformation Problem, producing over 1,900 diverse scenarios compatible with CARLA and LGSVL. We screen these in CARLA with Apollo to identify seed scenarios that trigger our target threat metrics.

% \subsection{Scenario Datasets for testing ADS}\label{sec:scenario_based_testing}
% Simulation-based testing has become the dominant paradigm for ADS evaluation. %Kalra and Paddock~\cite{kalra_driving_2016} demonstrated that proving the statistical safety of an autonomous vehicle would require over $10^9$ miles of on-road driving --- an impractical requirement that has driven the community towards virtual, scenario-based evaluation. High-fidelity simulators such as CARLA~\cite{dosovitskiy_carla_2017} and LGSVL~\cite{rong_lgsvl_2020} provide controllable virtual environments with realistic sensor models, physics, and rendering, but their effectiveness depends critically on the quality, diversity, and realism of the test scenarios they execute.
% Several naturalistic driving datasets have been developed to populate simulation testing pipelines. The inD dataset~\cite{bock_ind_2020} provides drone-recorded vehicle and pedestrian trajectories at German intersections, while highD~\cite{krajewski_highd_2018} captures highway traffic and rounD~\cite{krajewski_round_2020} records roundabout interactions. The Waymo Open Dataset~\cite{sun_scalability_2020} and nuScenes~\cite{caesar_nuscenes_2020} supply large-scale multi-sensor data from urban driving. The CARLA Leaderboard~\cite{noauthor_carla_nodate} defines ten traffic scenarios derived from the NHTSA pre-crash typology, providing a standardized benchmark for end-to-end ADS evaluation. The SCTrans dataset~\cite{dai_sctrans_2024}, which we use as the basis for our seed scenario identification, constructs over 1,900 simulation scenarios by transforming naturalistic traffic recordings into OpenSCENARIO files compatible with CARLA and LGSVL. SCTrans demonstrates that its scenarios can boost the collision-finding capability of existing ADS fuzzers by approximately seven times. However, a fundamental limitation shared by all these datasets is that they model \textit{well-behaved} traffic participants whose actions conform to expected driving norms, leaving adversarial edge cases underrepresented.

% \subsection{Simulation Platform and ADS Under Test}\label{sec:platform_ads}

% \textbf{CARLA} is an open-source driving simulator built on Unreal Engine that provides high-fidelity rendering, realistic vehicle physics, and flexible sensor models including camera, LiDAR, radar, and GPS~\cite{dosovitskiy_carla_2017}. It supports the OpenSCENARIO standard for scenario specification and the OpenDRIVE format for road network definition, enabling reproducible scenario execution. CARLA has been widely adopted in ADS testing research~\cite{kim_drivefuzz_2022, yan_-demand_2025, li_av-fuzzer_2020} and serves as the simulation platform for both our non-roundabout and roundabout experiments.

% \textbf{Apollo} is an open-source, industrial-grade L4 ADS that integrates perception (LiDAR and camera fusion), prediction, planning, and control modules into a full-stack driving pipeline~\cite{noauthor_apollo_nodate}. It has been extensively used as a test subject in ADS research~\cite{dai_sctrans_2024, kim_drivefuzz_2022, yan_-demand_2025, ji_effuzz_2023} due to its full-stack architecture and real-world deployment. Apollo's machine-learning-based pipeline represents the class of production-grade ADS most likely to be deployed in practice, making it a representative target for adversarial evaluation.

% \textbf{CARLA Traffic Manager (TM)} is the built-in rule-based autopilot provided by the CARLA simulator~\cite{dosovitskiy_carla_2017}. It controls vehicle behaviour through parameterizable rules governing target speed, following distance, lane-change probability, and traffic light compliance. Unlike Apollo's learning-based pipeline, TM follows deterministic, hand-crafted driving logic. We include it as a second ADS under test to compare the robustness of fundamentally different ADS architectures --- one machine-learning-based and one rule-based --- against identical adversarial scenarios.

% \textbf{SCTrans} is a transformation-based approach that constructs simulation scenario files from existing naturalistic driving datasets~\cite{dai_sctrans_2024}. Accepted at ICSE 2024, SCTrans formalises the conversion from traffic scenario recordings to OpenSCENARIO files as a Model Transformation Problem, producing a dataset of over 1,900 diverse scenarios that can directly test Apollo and Autoware via CARLA and LGSVL. We use 854 successfully simulated SCTrans scenarios as the starting point for our adversarial generation pipeline, screening them in CARLA with Apollo to identify seed scenarios that trigger our target threat metrics.

\subsection{NHTSA Pre-Crash Scenario Typology}\label{sec:NHTSA}

To ensure that the adversarial behaviours we generate reflect real-world collision patterns rather than arbitrary attack strategies, we ground our typology selection in the NHTSA Pre-Crash Scenario Typology for Crash Avoidance Research~\cite{noauthor_pre-crash_nodate}. This typology, widely adopted across the ADS community --- including by the CARLA Leaderboard~\cite{noauthor_carla_nodate} --- defines 37 pre-crash scenarios based on vehicle movements and critical events immediately prior to a crash, covering approximately 5.9 million police-reported light-vehicle crashes. We select ten two-vehicle typologies from this framework, grouped into three kinematic categories --- longitudinal, junction, and lateral --- that together account for approximately 71\% of all two-vehicle crashes. The selection criteria and category assignments are detailed in Sec.~\ref{sec:typology_selection}.

\subsection{Euro NCAP Assisted Driving Protocol}\label{sec:euro_ncap}

Having identified \textit{what} adversarial behaviours to generate, we require physically grounded thresholds to define \textit{when} a scenario becomes safety-critical. We derive these from the Euro NCAP Assisted Driving Highway \& Interurban Test and Assessment Protocol v2.2~\cite{noauthor_euro_nodate}, which specifies standardized test configurations for ACC, AEB, and Lane Support Systems across multiple collision types. The protocol's precisely specified vehicle speeds, following distances, and performance criteria provide an empirical basis for our two-tier threshold structure (critical and emergency) across five threat metrics. The specific threshold values and their derivations from the protocol are presented in Sec.~\ref{sec:threshold_selection} and Table~\ref{tab:thresholds}.